% Fuente: Pauta del Auxiliar 3, «Convexidad — Función convexa».
{\color{MidnightBlue}
\begin{enumerate}
  \item
\begin{align*}
    f(\lambda x + (1-\lambda) y)&=e^{g(\lambda x + (1-\lambda) y)}\\
    & \leq e^{\lambda g(x) + (1-\lambda) g(y)} \quad & \text{por convexidad de g(.)}\\
    & \leq \lambda e^{g(x)}+(1-\lambda)e^{g(y)} & \text{por convexidad de exp(.)}\\
    &= \lambda f(x) + (1-\lambda)f(y)
\end{align*}
  \item
Definiendo $g(x) =A x+Db$. Se puede ver que g(.) corresponde a una transformación lineal:
\begin{align*}
g(\lambda x + (1-\lambda)y)&=\lambda A x + (1-\lambda)A y + Db\\
&=\lambda(Ax+Db)+(1-\lambda)(Ay+Db)\\
&=\lambda g(x)+(1-\lambda)g(y)
\end{align*}

Luego
\begin{align*}
    f_r(\lambda x + (1-\lambda) y)&=f(g(\lambda x + (1-\lambda) y))\\
    & = f(\lambda g(x) + (1-\lambda) g(y)) \quad & \text{por linealidad de g(.)}\\
    & \leq \lambda f(g(x)) +(1-\lambda)f(g(y)) & \text{por convexidad de f(.)}\\
    &= \ \lambda f_r(x) +(1-\lambda)f_r(y)
\end{align*}
\end{enumerate}
}
